\documentclass[10pt,a4paper]{amsart}
\usepackage{amsmath,amssymb,amsthm,graphicx}
% Retained minimal preamble for the Stacks Project chapter sites-cohomology.tex (Cohomology on Sites).
% Verbatim excerpt of the chapter preamble: Xy-pic setup (native preamble lines 14-18) and the native
% theorem declarations with the mathematical macros (native preamble lines 154-217). Copied mechanically;
% nothing added, removed or redefined. The wrapper class and the amsmath/amssymb/amsthm/graphicx base
% packages are supplied by the assembler, so the class/externaldocument part of the native preamble is omitted.
% For commutative diagrams we use Xy-pic
\usepackage[all]{xy}

% We use 2cell for 2-commutative diagrams.
\xyoption{2cell}
% Theorem environments.
%
\theoremstyle{plain}
\newtheorem{theorem}[subsection]{Theorem}
\newtheorem{proposition}[subsection]{Proposition}
\newtheorem{lemma}[subsection]{Lemma}

\theoremstyle{definition}
\newtheorem{definition}[subsection]{Definition}
\newtheorem{example}[subsection]{Example}
\newtheorem{exercise}[subsection]{Exercise}
\newtheorem{situation}[subsection]{Situation}

\theoremstyle{remark}
\newtheorem{remark}[subsection]{Remark}
\newtheorem{remarks}[subsection]{Remarks}

\numberwithin{equation}{subsection}

% Macros
%
\def\lim{\mathop{\mathrm{lim}}\nolimits}
\def\colim{\mathop{\mathrm{colim}}\nolimits}
\def\Spec{\mathop{\mathrm{Spec}}}
\def\Hom{\mathop{\mathrm{Hom}}\nolimits}
\def\Ext{\mathop{\mathrm{Ext}}\nolimits}
\def\SheafHom{\mathop{\mathcal{H}\!\mathit{om}}\nolimits}
\def\SheafExt{\mathop{\mathcal{E}\!\mathit{xt}}\nolimits}
\def\Sch{\mathit{Sch}}
\def\Mor{\mathop{\mathrm{Mor}}\nolimits}
\def\Ob{\mathop{\mathrm{Ob}}\nolimits}
\def\Sh{\mathop{\mathit{Sh}}\nolimits}
\def\NL{\mathop{N\!L}\nolimits}
\def\CH{\mathop{\mathrm{CH}}\nolimits}
\def\proetale{{pro\text{-}\acute{e}tale}}
\def\etale{{\acute{e}tale}}
\def\QCoh{\mathit{QCoh}}
\def\Ker{\mathop{\mathrm{Ker}}}
\def\Im{\mathop{\mathrm{Im}}}
\def\Coker{\mathop{\mathrm{Coker}}}
\def\Coim{\mathop{\mathrm{Coim}}}

% Boxtimes
%
\DeclareMathSymbol{\boxtimes}{\mathbin}{AMSa}{"02}

%
% Macros for moduli stacks/spaces
%
\def\QCohstack{\mathcal{QC}\!\mathit{oh}}
\def\Cohstack{\mathcal{C}\!\mathit{oh}}
\def\Spacesstack{\mathcal{S}\!\mathit{paces}}
\def\Quotfunctor{\mathrm{Quot}}
\def\Hilbfunctor{\mathrm{Hilb}}
\def\Curvesstack{\mathcal{C}\!\mathit{urves}}
\def\Polarizedstack{\mathcal{P}\!\mathit{olarized}}
\def\Complexesstack{\mathcal{C}\!\mathit{omplexes}}
% \Pic is the operator that assigns to X its picard group, usage \Pic(X)
% \Picardstack_{X/B} denotes the Picard stack of X over B
% \Picardfunctor_{X/B} denotes the Picard functor of X over B
\def\Pic{\mathop{\mathrm{Pic}}\nolimits}
\def\Picardstack{\mathcal{P}\!\mathit{ic}}
\def\Picardfunctor{\mathrm{Pic}}
\def\Deformationcategory{\mathcal{D}\!\mathit{ef}}

% --- appended 2026-10-12 (Chief): the project own definitions of the reference/slogan/history
% display environments, copied mechanically from the shared Stacks preamble lines 8-12
% (\usepackage{verbatim} plus the three \newenvironment lines, which the chapter preamble keeps
% verbatim from the project preamble). Nothing added, removed or redefined: these environments
% carry no mathematical content and the project itself typesets them as comments. Without these
% three lines a counted statement whose own text uses them fails with Environment slogan undefined.
\usepackage{verbatim}
\newenvironment{reference}{\comment}{\endcomment}
\newenvironment{slogan}{\comment}{\endcomment}
\newenvironment{history}{\comment}{\endcomment}
\usepackage{cleveref}
\crefname{originallemma}{original lemma}{original lemmas}
\Crefname{originallemma}{Original lemma}{Original lemmas}
\makeatletter
\newlabel{derived-lemma-bounded-below-injectives-K-injective}{{070J}{source}{}{sourcepointer-1}{}}
\newlabel{derived-lemma-bounded-below-injectives-K-injective@cref}{{[originallemma][0][]070J}{[1][0][]source}{}{}{}}
\makeatother
\crefname{originallemma}{original lemma}{original lemmas}
\Crefname{originallemma}{Original lemma}{Original lemmas}
\makeatletter
\newlabel{homology-lemma-first-quadrant-ss}{{0132}{source}{}{sourcepointer-2}{}}
\newlabel{homology-lemma-first-quadrant-ss@cref}{{[originallemma][0][]0132}{[1][0][]source}{}{}{}}
\makeatother
\crefname{originalsection}{original section}{original sections}
\Crefname{originalsection}{Original section}{Original sections}
\makeatletter
\newlabel{section-internal-hom}{{017G}{source}{}{sourcepointer-3}{}}
\newlabel{section-internal-hom@cref}{{[originalsection][0][]017G}{[1][0][]source}{}{}{}}
\makeatother
\title{The internal hom of a bounded above complex of finite free summands with a bounded below complex is the termwise internal hom complex}
\author{The Stacks Project Authors}
\date{stacks-project:sites-cohomology.tex:local-snapshot-20261011; GNU FDL 1.2 or later}
\begin{document}
\maketitle
\noindent\textit{Editorial note.} All delivered mathematical text is the author's own, copied from the pinned Stacks Project chapter file sites-cohomology.tex; no statement or proof step was rewritten, reordered or removed.\par\medskip
\noindent\textit{Editorial note.} The counted proof is the author's complete argument: the proof is written out in full: the injective resolution, the finiteness of the products under the boundedness hypotheses, the identification of both complexes as total complexes of a double complex, and the spectral sequence argument reducing to the finite free case.\par\medskip
\noindent\textit{Editorial note.} The label lemma-Rhom-complex-of-direct-summands-finite-free names the counted statement, and the proof environment immediately following it in the source is delivered verbatim as its proof.\par\medskip
\noindent\textit{Editorial note.} The counted statement is carried by the author's own lemma environment in the chapter source.\par\medskip
\noindent\textit{Editorial note.} This card belongs to the priority queue of the D-72 twin ruling (the sites-side statement is the designated counting side). The statement and proof are copied from the pinned chapter file; the three cited inputs are kept as external pointers with their canonical Stacks tags.\par\medskip
\noindent\textit{Allowed standard prerequisites (editorial).}\begin{itemize}
\item Bounded below complexes of injectives are K-injective. Derived Categories, the lemma stating that a bounded below complex of injectives is K-injective. Supplies that the chosen resolution I\textasciicircum{}bullet is K-injective, so that the construction of the internal hom section applies.
\item The spectral sequence of a double complex. Homology, the lemma giving the first quadrant spectral sequence of a double complex. Supplies the two spectral sequences compared in the proof.
\item Construction of the internal hom complex. Same chapter, the section constructing the internal hom complex with its differential. Supplies the shape of the complexes H\textasciicircum{}bullet and (H')\textasciicircum{}bullet and the natural map between them.
\end{itemize}\medskip
\ifdefined\hypertarget\hypertarget{sourcepointer-1}{}\fi
\noindent\textit{Source reference.} Original lemma 070J: sources/stacks-chow/tags.txt, source lines 8958--8958. Prerequisite pointer invoked at proof line(s) 11984; it names a standard fact of a sibling chapter that the counted argument uses, and every step of the counted proof built on it is retained verbatim in the delivered body.\par\medskip
\ifdefined\hypertarget\hypertarget{sourcepointer-2}{}\fi
\noindent\textit{Source reference.} Original lemma 0132: sources/stacks-chow/tags.txt, source lines 1420--1420. Prerequisite pointer invoked at proof line(s) 12005; it names a standard fact of a sibling chapter that the counted argument uses, and every step of the counted proof built on it is retained verbatim in the delivered body.\par\medskip
\ifdefined\hypertarget\hypertarget{sourcepointer-3}{}\fi
\noindent\textit{Source reference.} Original section 017G: sources/stacks-sites-cohomology/pinned-original-source.tex, source lines 8836--8836. Prerequisite pointer invoked at proof line(s) 11976; it names a standard fact of the same chapter that the counted argument uses, and every step of the counted proof built on it is retained verbatim in the delivered body.\par\medskip
\begin{lemma}
\label{lemma-Rhom-complex-of-direct-summands-finite-free}
Let $(\mathcal{C}, \mathcal{O})$ be a ringed site.
Let $\mathcal{E}^\bullet$, $\mathcal{F}^\bullet$ be complexes
of $\mathcal{O}$-modules with
\begin{enumerate}
\item $\mathcal{F}^n = 0$ for $n \ll 0$,
\item $\mathcal{E}^n = 0$ for $n \gg 0$, and
\item $\mathcal{E}^n$ isomorphic to a direct summand of a finite
free $\mathcal{O}$-module.
\end{enumerate}
Then the internal hom $R\SheafHom(\mathcal{E}^\bullet, \mathcal{F}^\bullet)$
is represented by the complex $\mathcal{H}^\bullet$ with terms
$$
\mathcal{H}^n =
\bigoplus\nolimits_{n = p + q}
\SheafHom_\mathcal{O}(\mathcal{E}^{-q}, \mathcal{F}^p)
$$
and differential as described in Section \ref{section-internal-hom}.
\end{lemma}

\begin{proof}
Choose a quasi-isomorphism $\mathcal{F}^\bullet \to \mathcal{I}^\bullet$
where $\mathcal{I}^\bullet$ is a bounded below complex of injectives.
Note that $\mathcal{I}^\bullet$ is K-injective
(Derived Categories, Lemma
\ref{derived-lemma-bounded-below-injectives-K-injective}).
Hence the construction in Section \ref{section-internal-hom}
shows that
$R\SheafHom(\mathcal{E}^\bullet, \mathcal{F}^\bullet)$ is 
represented by the complex $(\mathcal{H}')^\bullet$ with terms
$$
(\mathcal{H}')^n =
\prod\nolimits_{n = p + q}
\SheafHom_\mathcal{O}(\mathcal{E}^{-q}, \mathcal{I}^p) =
\bigoplus\nolimits_{n = p + q}
\SheafHom_\mathcal{O}(\mathcal{E}^{-q}, \mathcal{I}^p)
$$
(equality because there are only finitely many nonzero terms).
Note that $\mathcal{H}^\bullet$ is the total complex associated to
the double complex with terms
$\SheafHom_\mathcal{O}(\mathcal{E}^{-q}, \mathcal{F}^p)$
and similarly for $(\mathcal{H}')^\bullet$.
The natural map $(\mathcal{H}')^\bullet \to \mathcal{H}^\bullet$
comes from a map of double complexes.
Thus to show this map is a quasi-isomorphism, we may use the spectral
sequence of a double complex
(Homology, Lemma \ref{homology-lemma-first-quadrant-ss})
$$
{}'E_1^{p, q} =
H^p(\SheafHom_\mathcal{O}(\mathcal{E}^{-q}, \mathcal{F}^\bullet))
$$
converging to $H^{p + q}(\mathcal{H}^\bullet)$ and similarly for
$(\mathcal{H}')^\bullet$. To finish the proof of the lemma it
suffices to show that $\mathcal{F}^\bullet \to \mathcal{I}^\bullet$
induces an isomorphism
$$
H^p(\SheafHom_\mathcal{O}(\mathcal{E}, \mathcal{F}^\bullet))
\longrightarrow
H^p(\SheafHom_\mathcal{O}(\mathcal{E}, \mathcal{I}^\bullet))
$$
on cohomology sheaves whenever $\mathcal{E}$ is a direct summand of a
finite free $\mathcal{O}$-module. Since this is clear when $\mathcal{E}$
is finite free the result follows.
\end{proof}

\end{document}
